\documentclass[11pt]{article}
\newcommand{\DocNumber}{VOLUME 03}
\newcommand{\DocTitle}{Modern C++23 Programming Handbook}
\newcommand{\DocSubtitle}{Applied patterns and hands-on implementation}
\newcommand{\DocSummary}{Twelve implementable engineering patterns with design contracts, original samples and an integrated capstone lab.}
% Shared, pdfLaTeX-only design system for the C++23 documentation suite.
\usepackage[T1]{fontenc}
\usepackage[utf8]{inputenc}
\usepackage{lmodern}
\usepackage[english]{babel}
\usepackage[letterpaper,margin=0.82in,headheight=16pt,headsep=17pt]{geometry}
\usepackage{microtype}
\usepackage{amsmath,amssymb}
\usepackage{array,booktabs,tabularx,longtable}
\usepackage[table]{xcolor}
\usepackage[most]{tcolorbox}
\usepackage{listings}
\usepackage{fancyhdr}
\usepackage{titlesec}
\usepackage{enumitem}
\usepackage{hyperref}
\usepackage{needspace}
\usepackage{etoolbox}
\definecolor{ink}{HTML}{13253B}
\definecolor{teal}{HTML}{006E7F}
\definecolor{muted}{HTML}{5A6C7C}
\definecolor{fog}{HTML}{EDF4F6}
\definecolor{sand}{HTML}{FAF3E8}
\definecolor{rule}{HTML}{C9D7DF}
\definecolor{codebg}{HTML}{F7F9FB}
\definecolor{codecomment}{HTML}{47715D}
\hypersetup{colorlinks=true,linkcolor=teal,urlcolor=teal,pdfauthor={Independent C++23 engineering documentation},pdftitle={C++23 Technical Suite}}
\setlength{\parindent}{0pt}
\setlength{\parskip}{6pt}
\setcounter{tocdepth}{2}
\renewcommand{\arraystretch}{1.18}
\titleformat{\section}{\Large\bfseries\color{ink}}{\thesection}{0.7em}{}
\titleformat{\subsection}{\large\bfseries\color{teal}}{\thesubsection}{0.55em}{}
\titlespacing*{\section}{0pt}{17pt}{7pt}
\titlespacing*{\subsection}{0pt}{11pt}{4pt}
\pagestyle{fancy}
\fancyhf{}
\fancyhead[L]{\footnotesize\bfseries\color{teal}C++23 / ISO 14882:2024}
\fancyhead[R]{\footnotesize\color{muted}\DocNumber}
\fancyfoot[L]{\scriptsize\color{muted}Independent engineering companion, not an ISO publication}
\fancyfoot[R]{\scriptsize\color{muted}\thepage}
\renewcommand{\headrulewidth}{0.35pt}
\lstdefinestyle{cppstyle}{language=C++,basicstyle=\ttfamily\footnotesize,keywordstyle=\color{teal}\bfseries,commentstyle=\color{codecomment},stringstyle=\color{ink},backgroundcolor=\color{codebg},frame=single,rulecolor=\color{rule},breaklines=true,breakatwhitespace=true,columns=fullflexible,keepspaces=true,showstringspaces=false,tabsize=2,numbers=none,upquote=true,xleftmargin=3pt,xrightmargin=3pt,aboveskip=6pt,belowskip=8pt}
\lstset{style=cppstyle}
\newtcolorbox{standardbox}{enhanced,breakable,colback=fog,colframe=teal!35!white,boxrule=0.55pt,arc=1.4mm,left=3mm,right=3mm,top=2mm,bottom=2mm,before skip=7pt,after skip=8pt}
\newtcolorbox{practicebox}{enhanced,breakable,colback=sand,colframe=ink!20!white,boxrule=0.55pt,arc=1.4mm,left=3mm,right=3mm,top=2mm,bottom=2mm,before skip=7pt,after skip=8pt}
\newcommand{\anchorref}[1]{\par\smallskip{\footnotesize\textbf{Standard locator:} \texttt{#1}}\par}
\newcommand{\standard}[1]{\begin{standardbox}\textbf{Standard-grounded interpretation.} #1\end{standardbox}}
\newcommand{\practice}[1]{\begin{practicebox}\textbf{Engineering recommendation (not a normative ISO rule).} #1\end{practicebox}}
\newcommand{\takeaway}[1]{\par\textbf{Review question.} #1\par}
\newcommand{\docfront}{%
\begin{titlepage}
\vspace*{25mm}
{\color{teal}\bfseries\large ENGINEERING REFERENCE SERIES\par}
\vspace{12mm}
{\raggedright\color{ink}\bfseries\fontsize{24}{30}\selectfont\DocTitle\par}
\vspace{7mm}
{\color{muted}\fontsize{15}{19}\selectfont\DocSubtitle\par}
\vspace{12mm}
\noindent\color{teal}\rule{\linewidth}{2pt}\par
\vspace{9mm}
{\large\DocSummary\par}
\vfill
\begin{tabularx}{\linewidth}{@{}lX@{}}
\textbf{Volume} & \DocNumber\\
\textbf{Standards basis} & ISO/IEC 14882:2024 (seventh edition, C++23)\\
\textbf{Document type} & Original interpretive and instructional material\\
\textbf{Date} & 16 October 2026\\
\textbf{Toolchain} & pdfLaTeX; illustrative C++23 examples require separate compiler support\\
\end{tabularx}
\vspace{8mm}
{\footnotesize This document is not an authorized copy of, replacement for, or amendment to ISO/IEC 14882:2024. The original standard controls in the event of disagreement.}\par
\end{titlepage}
\pagenumbering{roman}
\tableofcontents
\clearpage
\pagenumbering{arabic}
\section*{How to use this volume}\addcontentsline{toc}{section}{How to use this volume}
The reference is \emph{ISO/IEC 14882:2024, Programming languages --- C++}. Section locators and square-bracketed stable names identify where to verify precise wording; they are not quotations. \textbf{Standard-grounded interpretation} boxes summarize language or library semantics. \textbf{Engineering recommendation} boxes contain optional design, security, portability, or process guidance. Code is original and illustrative; individual library facilities may not yet be implemented by every vendor. This suite does not grant rights to reproduce or distribute the licensed standard.\par\medskip
}
\newcommand{\docend}{\clearpage\section*{Reference and traceability}\addcontentsline{toc}{section}{Reference and traceability}
\textbf{Primary source.} ISO/IEC 14882:2024(en), \emph{Programming languages --- C++}, seventh edition, October 2024. Use its numbered clauses, stable names, Annex C (informative compatibility), and Annex D (compatibility features) for authoritative details.\par
\textbf{Limits.} Clause-level summaries and exercises in this document are selective explanations, not exhaustive reproductions of all semantic constraints, library mandates, complexity clauses, and overload participation requirements. A successful compiler test demonstrates behavior of that specific implementation and configuration, not universal conformance.\par
\textbf{Use of contemporary engineering tools.} Compilers, sanitizers, analyzers, CI systems, repository policies, and vulnerability taxonomies are implementation guidance rather than language-standard requirements.\par}

\begin{document}
\docfront
\section{Implementation workflow}
For each pattern, implement the smallest invariant-preserving change, compile with the chosen C++23 toolchain, test its negative cases, run enabled diagnostics and measure performance only after correctness is established.

\small
\begin{longtable}{@{}p{.18\linewidth}p{.40\linewidth}p{.34\linewidth}@{}}
\toprule
\textbf{Stage} & \textbf{Artifact} & \textbf{Exit condition}\\
\midrule
\endfirsthead
\toprule
\textbf{Stage} & \textbf{Artifact} & \textbf{Exit condition}\\
\midrule
\endhead
Contract & Interface and invariants & Valid states and error outcomes described\\
\addlinespace[2pt]
Prototype & Compiled C++23 implementation & No unreviewed extensions\\
\addlinespace[2pt]
Verification & Unit, negative and sanitizer tests & Expected outcomes reproducible\\
\addlinespace[2pt]
Integration & Dependency, ABI and build notes & Cross-module behavior evaluated\\
\addlinespace[2pt]
Operations & Metrics and rollback path & Regression detection and reversibility\\
\addlinespace[2pt]
\bottomrule
\end{longtable}
\normalsize
\section{Project baseline and dependency architecture}
\anchorref{4--5; 10; 16}
Establish a build that makes language mode, target triple, standard library, warnings, and dependency versions observable. A standards-conforming program and a reproducible build solve different problems.

\standard{The cited clauses establish the underlying C++ semantics used by this implementation. The higher-level architecture pattern is an engineering interpretation, not an ISO-mandated design.}
\begin{itemize}[leftmargin=*,itemsep=1pt,topsep=3pt]
\item Define preconditions, postconditions and state ownership.
\item Write a positive test and a meaningful negative/boundary test.
\item Review exception and cancellation paths.
\end{itemize}
\begin{lstlisting}[language=C++]
cmake_minimum_required(VERSION 3.25)
project(engineering CXX)
add_executable(example main.cpp)
set_target_properties(example PROPERTIES
    CXX_STANDARD 23 CXX_STANDARD_REQUIRED YES
    CXX_EXTENSIONS NO)
# Probe vendor library features separately.
\end{lstlisting}
\practice{Lock toolchains for CI; check the selected implementation rather than assuming full C++23 coverage.}
\takeaway{Hands-on task: Create a manifest describing compiler, library, flags and verified feature macros.}
\section{Value types with explicit invariants}
\anchorref{9; 11}
Express valid states through constructors and private representation. Value semantics simplify test isolation, caching, and concurrency when mutable shared state is absent.

\standard{The cited clauses establish the underlying C++ semantics used by this implementation. The higher-level architecture pattern is an engineering interpretation, not an ISO-mandated design.}
\begin{itemize}[leftmargin=*,itemsep=1pt,topsep=3pt]
\item Define preconditions, postconditions and state ownership.
\item Write a positive test and a meaningful negative/boundary test.
\item Review exception and cancellation paths.
\end{itemize}
\begin{lstlisting}[language=C++]
#include <expected>
#include <string_view>
struct Port {
private:
    unsigned short value_;
    explicit Port(unsigned short value) : value_(value) {}
public:
    static std::expected<Port, std::string_view> make(int value) {
        if (value < 1 || value > 65535) return
            std::unexpected("port out of range");
        return Port{static_cast<unsigned short>(value)};
    }
    unsigned short value() const noexcept { return value_; }
};
\end{lstlisting}
\practice{Use named factory functions when construction can fail; prevent invalid states rather than relying on documentation.}
\takeaway{Hands-on task: Prove that no public operation can construct an invalid port.}
\section{RAII and linear resource release}
\anchorref{6.7; 11; 20; 33}
Resource acquisition is initialization ties resource release to object lifetime and structured destruction; it is a design idiom enabled by core lifetime rules.

\standard{The cited clauses establish the underlying C++ semantics used by this implementation. The higher-level architecture pattern is an engineering interpretation, not an ISO-mandated design.}
\begin{itemize}[leftmargin=*,itemsep=1pt,topsep=3pt]
\item Define preconditions, postconditions and state ownership.
\item Write a positive test and a meaningful negative/boundary test.
\item Review exception and cancellation paths.
\end{itemize}
\begin{lstlisting}[language=C++]
#include <mutex>
class Counter {
    mutable std::mutex mu_;
    int count_{};
public:
    void increment() {
        std::scoped_lock lock{mu_};
        ++count_;
    }
    int get() const {
        std::scoped_lock lock{mu_};
        return count_;
    }
};
\end{lstlisting}
\practice{Use dedicated guard objects even when the protected function currently cannot throw.}
\takeaway{Hands-on task: Design a failing operation and demonstrate that the lock is released.}
\section{Borrowing, ownership, and non-owning views}
\anchorref{6.7; 20; 23; 24}
Borrowing allows efficient APIs without ownership transfer, but a view can outlive its storage. Document the temporal contract.

\standard{The cited clauses establish the underlying C++ semantics used by this implementation. The higher-level architecture pattern is an engineering interpretation, not an ISO-mandated design.}
\begin{itemize}[leftmargin=*,itemsep=1pt,topsep=3pt]
\item Define preconditions, postconditions and state ownership.
\item Write a positive test and a meaningful negative/boundary test.
\item Review exception and cancellation paths.
\end{itemize}
\begin{lstlisting}[language=C++]
#include <string>
#include <string_view>
std::string add_prefix(std::string_view input) {
    std::string result{"id:"};
    result.append(input);
    return result; // owns the returned bytes
}
\end{lstlisting}
\practice{Return owned values across asynchronous, persistent or remote boundaries.}
\takeaway{Hands-on task: Write a compile-and-run regression test for dangling views with sanitizers.}
\section{Moving values and the rule of zero}
\anchorref{11.4; 20}
Avoid manually writing destructor/copy/move members unless the class itself directly manages a non-RAII resource. When move is added, reason about the source and destination states.

\standard{The cited clauses establish the underlying C++ semantics used by this implementation. The higher-level architecture pattern is an engineering interpretation, not an ISO-mandated design.}
\begin{itemize}[leftmargin=*,itemsep=1pt,topsep=3pt]
\item Define preconditions, postconditions and state ownership.
\item Write a positive test and a meaningful negative/boundary test.
\item Review exception and cancellation paths.
\end{itemize}
\begin{lstlisting}[language=C++]
#include <string>
#include <vector>
struct Job {
    std::string id;
    std::vector<int> samples;
};
Job make_job() { return {"abc", {1, 2, 3}}; }
\end{lstlisting}
\practice{Prefer composition with vector, string and smart pointers so generated special members do the right thing.}
\takeaway{Hands-on task: Test copy, move, self-assignment (when supported) and exception paths.}
\section{Explicit interfaces with concepts}
\anchorref{13; 18}
Generic code should describe the operations and semantics it requires. Constraints improve diagnosis and overload participation but are not full formal proofs.

\standard{The cited clauses establish the underlying C++ semantics used by this implementation. The higher-level architecture pattern is an engineering interpretation, not an ISO-mandated design.}
\begin{itemize}[leftmargin=*,itemsep=1pt,topsep=3pt]
\item Define preconditions, postconditions and state ownership.
\item Write a positive test and a meaningful negative/boundary test.
\item Review exception and cancellation paths.
\end{itemize}
\begin{lstlisting}[language=C++]
#include <concepts>
template<typename T>
concept Summable = requires(T a, T b) {
    { a + b } -> std::same_as<T>;
};
template<Summable T>
T add(T a, T b) { return a + b; }
\end{lstlisting}
\practice{Constrain APIs at boundaries; avoid unrelated template instantiation explosions.}
\takeaway{Hands-on task: Provide both accepted and intentionally rejected types as tests.}
\section{Error-oriented API design with expected}
\anchorref{14; 19; 22.8}
Use expected when a computation naturally returns either a value or a recoverable error. This does not replace exceptions automatically; separate expected errors from programmer contract violations.

\standard{The cited clauses establish the underlying C++ semantics used by this implementation. The higher-level architecture pattern is an engineering interpretation, not an ISO-mandated design.}
\begin{itemize}[leftmargin=*,itemsep=1pt,topsep=3pt]
\item Define preconditions, postconditions and state ownership.
\item Write a positive test and a meaningful negative/boundary test.
\item Review exception and cancellation paths.
\end{itemize}
\begin{lstlisting}[language=C++]
#include <expected>
enum class ParseError { empty, bad_digit };
std::expected<int, ParseError> digit(char c) {
    if (c < '0' || c > '9')
        return std::unexpected(ParseError::bad_digit);
    return c - '0';
}
\end{lstlisting}
\practice{Design stable, structured error categories rather than string matching.}
\takeaway{Hands-on task: Specify how caller code propagates and records each error.}
\section{Range pipelines and complexity budgets}
\anchorref{24--27}
Lazy pipelines can cut allocations but may recompute projections or predicates. Iterator category, range ownership and complexity still matter.

\standard{The cited clauses establish the underlying C++ semantics used by this implementation. The higher-level architecture pattern is an engineering interpretation, not an ISO-mandated design.}
\begin{itemize}[leftmargin=*,itemsep=1pt,topsep=3pt]
\item Define preconditions, postconditions and state ownership.
\item Write a positive test and a meaningful negative/boundary test.
\item Review exception and cancellation paths.
\end{itemize}
\begin{lstlisting}[language=C++]
#include <ranges>
#include <vector>
std::vector<int> values{2, 7, 4, 9};
auto even_doubled = values
    | std::views::filter([](int x){ return x % 2 == 0; })
    | std::views::transform([](int x){ return x * 2; });
\end{lstlisting}
\practice{Benchmark on representative input and inspect evaluation frequency.}
\takeaway{Hands-on task: Count predicate invocations under repeated traversal.}
\section{Module boundaries and distribution}
\anchorref{10; 15--16}
Modules introduce new build dependencies and export/reachability boundaries. A migration should prioritize stable APIs and minimal macro interaction.

\standard{The cited clauses establish the underlying C++ semantics used by this implementation. The higher-level architecture pattern is an engineering interpretation, not an ISO-mandated design.}
\begin{itemize}[leftmargin=*,itemsep=1pt,topsep=3pt]
\item Define preconditions, postconditions and state ownership.
\item Write a positive test and a meaningful negative/boundary test.
\item Review exception and cancellation paths.
\end{itemize}
\begin{lstlisting}[language=C++]
// geometry.cppm (interface unit)
export module geometry;
export struct Point { double x{}, y{}; };
export double length2(Point p) {
    return p.x * p.x + p.y * p.y;
}
\end{lstlisting}
\practice{Maintain a header-based compatibility path if required by your target compilers.}
\takeaway{Hands-on task: Draw import dependencies and identify build-order cycles.}
\section{Concurrency ownership and cancellation}
\anchorref{6.9; 33}
Safe parallel design needs both synchronization and lifetime control. Relaxed atomics prevent data races on the atomic itself but do not synchronize other objects.

\standard{The cited clauses establish the underlying C++ semantics used by this implementation. The higher-level architecture pattern is an engineering interpretation, not an ISO-mandated design.}
\begin{itemize}[leftmargin=*,itemsep=1pt,topsep=3pt]
\item Define preconditions, postconditions and state ownership.
\item Write a positive test and a meaningful negative/boundary test.
\item Review exception and cancellation paths.
\end{itemize}
\begin{lstlisting}[language=C++]
#include <stop_token>
#include <thread>
#include <atomic>
void run(std::stop_token stop, std::atomic<int>& count) {
    while (!stop.stop_requested()) {
        count.fetch_add(1, std::memory_order_relaxed);
        break; // example does bounded work
    }
}
\end{lstlisting}
\practice{Prefer message passing or immutable data where feasible, and choose jthread when cooperative stop ownership is useful.}
\takeaway{Hands-on task: Describe the happens-before edges and shutdown order.}
\section{Testable IO boundary and dependency injection}
\anchorref{19; 31}
Separate parsing, business rules and infrastructure adapters so deterministic logic can be tested without filesystem or network state.

\standard{The cited clauses establish the underlying C++ semantics used by this implementation. The higher-level architecture pattern is an engineering interpretation, not an ISO-mandated design.}
\begin{itemize}[leftmargin=*,itemsep=1pt,topsep=3pt]
\item Define preconditions, postconditions and state ownership.
\item Write a positive test and a meaningful negative/boundary test.
\item Review exception and cancellation paths.
\end{itemize}
\begin{lstlisting}[language=C++]
#include <istream>
#include <expected>
enum class ReadError { unavailable, invalid };
std::expected<int, ReadError> read_number(std::istream& input) {
    int value{};
    if (!(input >> value)) return
        std::unexpected(ReadError::invalid);
    return value;
}
\end{lstlisting}
\practice{Represent IO failure distinctly from malformed input, and make effectful dependencies explicit.}
\takeaway{Hands-on task: Inject malformed and truncated input and assert exact error codes.}
\section{Performance and allocation discipline}
\anchorref{20; 24--27}
Complexity promises are contractual but a fast asymptotic algorithm is not automatically fastest for a given cache or data distribution.

\standard{The cited clauses establish the underlying C++ semantics used by this implementation. The higher-level architecture pattern is an engineering interpretation, not an ISO-mandated design.}
\begin{itemize}[leftmargin=*,itemsep=1pt,topsep=3pt]
\item Define preconditions, postconditions and state ownership.
\item Write a positive test and a meaningful negative/boundary test.
\item Review exception and cancellation paths.
\end{itemize}
\begin{lstlisting}[language=C++]
#include <vector>
void build_ids(std::vector<int>& ids, int n) {
    ids.clear();
    if (n > 0) ids.reserve(static_cast<std::size_t>(n));
    for (int i = 0; i < n; ++i) ids.push_back(i);
}
\end{lstlisting}
\practice{Start with a workload model, optimize only after measurement, and retain a correctness oracle.}
\takeaway{Hands-on task: Measure allocations and p95 latency separately from total throughput.}
\section{Integrated lab: bounded job processor}
Build a library that accepts job records, validates their identifiers and sizes, processes a range of jobs, and returns per-job expected results. Keep parsing, core transformation and output separate. Extend the implementation to run jobs with jthread and stop\_token only after the serial version is verified.

\begin{itemize}[leftmargin=*,itemsep=1pt,topsep=3pt]
\item Input contract: bounded record count, bounded sizes, explicit encoding, reject malformed jobs.
\item Type contract: opaque validated IDs; no raw owning pointer in the public API.
\item Failure contract: structured expected error for invalid input; exception-safety guarantee for allocation failures.
\item Concurrency contract: define queue ownership and safe stopping before adding threads.
\item Verification: property tests, AddressSanitizer/UndefinedBehaviorSanitizer, ThreadSanitizer in separate builds, and a benchmark.
\end{itemize}
\docend
\end{document}
